\documentclass[11pt]{article}
\usepackage[margin=0.95in]{geometry}
\usepackage{lmodern}
\usepackage[T1]{fontenc}
\usepackage{microtype}
\usepackage{xcolor}
\usepackage{amsmath,amssymb,mathtools}
\usepackage{booktabs,tabularx,longtable,array,multirow}
\newcolumntype{Y}{>{\raggedright\arraybackslash}X}
\usepackage{enumitem}
\usepackage{caption}
\usepackage{listings}
\usepackage{tikz}
\usetikzlibrary{arrows.meta,positioning,shapes.geometric,calc}
\usepackage[hidelinks]{hyperref}

\definecolor{navy}{HTML}{1B3A4B}
\definecolor{teal}{HTML}{0E6B66}
\definecolor{tealfaint}{HTML}{E5F3F1}
\definecolor{gold}{HTML}{A56B10}
\definecolor{goldfaint}{HTML}{FBF3E0}
\definecolor{brick}{HTML}{8E3A3A}
\definecolor{brickfaint}{HTML}{F8E8E6}
\definecolor{leaf}{HTML}{2F6F4E}
\definecolor{leaffaint}{HTML}{E5F2EA}
\definecolor{slate}{HTML}{5E6A75}
\definecolor{linec}{HTML}{243140}
\definecolor{paper}{HTML}{F7F4EF}
\definecolor{unk}{HTML}{E4E1DA}
\definecolor{freec}{HTML}{F7F7F5}
\definecolor{occ}{HTML}{2C333A}
\definecolor{enc}{HTML}{C45C4A}
\definecolor{frcell}{HTML}{2A9D8F}
\definecolor{codebg}{HTML}{F4F1EA}

\lstset{
  basicstyle=\ttfamily\footnotesize,
  backgroundcolor=\color{codebg},
  frame=single,
  rulecolor=\color{linec},
  breaklines=true,
  columns=fullflexible,
  keepspaces=true,
  xleftmargin=1em,
  framexleftmargin=0.6em
}

\setlist[itemize]{leftmargin=1.2em, itemsep=0.2em, topsep=0.35em}
\setlist[enumerate]{leftmargin=1.4em, itemsep=0.25em, topsep=0.35em}
\captionsetup{font=small, labelfont=bf, skip=6pt}
\setlength{\parskip}{0.35em}
\setlength{\parindent}{1.15em}
\setlength{\emergencystretch}{2em}
\renewcommand{\arraystretch}{1.15}

\newcommand{\code}[1]{\texttt{#1}}
\newcommand{\commit}[1]{\texttt{#1}}


\title{Publication and Intellectual-Property Assessment\\of the Frontier Exploration Policy}
\author{Internal review. Not legal advice.\\Not a patent opinion.}
\date{Based on \texttt{main} at \texttt{de8209a}, explorer at \texttt{eccbb97}.\\Public repository: \texttt{amey1942007/Ros\_lidar\_bot}.}

\begin{document}
\maketitle

\begin{abstract}
\noindent
The exploration code is a classical frontier detector with a failure-management policy. That policy is specific, documented, and internally consistent enough to describe. It is not supported by a controlled experiment. A leading robotics conference will not accept a performance claim in the current state, and a patent filing is a poor fit for a method that is already public under Apache-2.0. The highest-value next step is a logged comparison against nearest-frontier and a path-cost utility, after three small correctness fixes. Counsel should be asked about the disclosure clock before any of this text is posted outside the group.
\end{abstract}

\section{Judgment}
Two different questions are in play, and they have different answers.

\textbf{Can this be submitted to IROS, ICRA, or RA-L as it stands?} No, not as a paper that claims the robot explores better. The method section of such a paper could be written from the companion report. The results section cannot, because the results do not exist. A submission that invents them, or that treats commit messages as trials, would be rejected and would be wrong.

\textbf{Is the method worth protecting as intellectual property?} Unlikely in its present form, and possibly already given away. The repository is public. The licence is Apache-2.0, which includes a patent grant from contributors to users of the code. The algorithmic core is prior art. What remains is a combination of heuristics. Combinations can be patented in some countries, but only if they are new, non-obvious, and not already disclosed. This assessment is not a clearance to file, and it is not a statement that a filing is impossible. It is a statement that a filing made without a lawyer, after this public history, is the wrong default.

\section{What the work actually is}
The detector is Yamauchi's idea, implemented as unknown cells beside free cells, clustered on the full grid, and scored by a weighted sum. \texttt{explore\_lite} already scores a gain against a cost and blacklists failures. Information-gain explorers already do something stricter than cluster size. Fast frontier detectors already avoid the full-grid scan this code still performs.

The local packaging, which is the only part that should be discussed as the authors' work, is:
\begin{enumerate}
\item a border-seeded flood fill that removes enclosed unknown cells and stores them as rounded world keys rather than a disk;
\item a counter that treats Nav2 \emph{success} with residual unknown as different from an abort, and blacklists at two such arrivals;
\item a standby queue of separated frontiers that a safety-stop Boolean can preempt, with a clean cancel when the frontier disappears during the trip;
\item a two-stage abort rule whose stated purpose is to ignore one-off failures while the map is growing;
\item a sticky relaxation schedule tied to a one-shot reading of the first grid.
\end{enumerate}
Items 4 and 5 are ordinary engineering. Items 1--3 are the combination a reviewer or an examiner would actually look at. None of them has a measured effect.

The design history is a point in the authors' favour if it is told honestly. A $0.25$ openness floor and a $0.60$\,m border margin were removed because, the commits say, they rejected real frontiers. That is evidence of engineering judgment. It is not evidence of a better explorer.

\section{Conference readiness}
\subsection{What a programme committee will look for}
IROS and ICRA systems papers can be accepted without a new theorem. They are not accepted without a comparison. The usual bar for an exploration paper is: a stated baseline, more than one environment, repeated trials, and an ablation of the piece the authors say matters. RA-L is at least that strict, and it expects the letter to stand alone.

Against that bar the repository currently provides the method and withholds the experiment. Commit notes that a transform tolerance or a safety cone was ``verified live'' are anecdotes about navigation, not exploration trials. They have no sample size.

\subsection{How the contribution would be scored today}
\begin{table}[ht]
\centering
\caption{Readiness against the questions a reviewer will ask. Scores are judgments of the present record, not of a future paper.}
\small
\begin{tabular}{@{}p{4.2cm}p{2.2cm}p{7.6cm}@{}}
\toprule
Question & State & Why \\
\midrule
Is the problem real? & Adequate & Frontiers that cannot be observed, and safety stops that pin the robot, are real failure modes. \\
Is the method specified? & Adequate, after this report & The June 2026 note is not a specification. The source is. \\
Is it distinguished from prior work? & Weak & The text must concede Yamauchi, utility scoring, and \texttt{explore\_lite}. The residual claim is the failure policy. \\
Is there evidence? & Missing & No bags, no tables, no baselines. \\
Are limitations disclosed? & Now yes, in the two manuscripts & They were not disclosed in the June note, which also contradicts the code. \\
Can a stranger reproduce the configuration? & Partial & Launch files and YAML are public. Worlds used in the May commits are not on \texttt{main}. Dependency versions are not pinned. \\
\bottomrule
\end{tabular}
\end{table}

\subsection{What would make a submission defensible}
In order of impact:

\begin{enumerate}
\item \textbf{Run the comparison.} Nearest-frontier and a size-over-path-cost utility, same detector, same Nav2 file, at least twenty repeats, office and warehouse if those worlds can be recovered, plus one hardware floor. Outcomes: time and path to a pre-declared coverage level, repeated goals within $1.2$\,m, safety preemptions, aborts, phantom counts. Paired tests. This single item changes the paper from a description into a result, whether the result is positive or not.
\item \textbf{Ablate the two mechanisms that are actually argued.} Enclosed-cell rejection, and the phantom counter. If they do not move the repeated-goal count, do not claim them. The queue and the far bonus are secondary ablations.
\item \textbf{Repair the score before measuring it.} A popped goal currently returns $-\mathrm{distance}/d_{\max}$, not the four-term utility. A comparison that logs ``score'' will mix two functions. The border test should be a distance; \texttt{int(0.15/0.05)} is $2$, not the $3$ cells the comments describe. The completion log should latch or be deleted, so a trial does not record a finish that the node does not perform.
\item \textbf{Instrument goals.} One structured line per goal: time, pose, both scores if both exist, cluster size, end reason, whether the queue was used, unknown cells inside $0.8$\,m. The existing debug log is not that record.
\item \textbf{Write the paper as a failure-policy paper, not as a new explorer.} The title should not say ``efficient'' or ``optimal''. The related-work section should put \texttt{explore\_lite} and Holz \emph{et al.} in the first page of that section. The June report should not be cited as prior experimental evidence.
\item \textbf{Only then choose the venue.} A positive, ablated result on coverage and repeated failures can aim at IROS or ICRA. A negative result, if the protocol was pre-declared, is also publishable and is more honest than a tuned demo. A workshop or a late-breaking abstract is the right place if the trials are small. RA-L should wait until the result is stable and the letter is short.
\end{enumerate}

A paper submitted before item 1 is a method note with no method-note venue that will treat uncompared heuristics as a contribution. Do not submit it to the main tracks in order to ``see what reviewers say.''

\subsection{Defects that a reviewer will find if the code is released with the paper}
These are already visible statically. They should be fixed or listed as limitations in the submission itself.
\begin{itemize}
\item Pop score and selection score are different functions, and idle hysteresis compares them.
\item The idle loop prefers a queue of runners-up for up to $20$\,s and can dispatch beyond the normal search radius, up to the $12$\,m cap.
\item Permanent-zone keys alias under Python's half-to-even \texttt{round}. The RViz disk is larger than the predicate.
\item The phantom path can store $1200$\,s while logging ``10 min'', because it temporarily changes the timeout base before the exponential update.
\item A safety preemption can call \texttt{cancel} on a goal handle that is still \texttt{None}.
\item The safety Boolean treats a dead wheel encoder as a reason to blacklist a goal.
\item The map subscription uses default volatile QoS. Whether that drops the first SLAM map was not tested. It should be tested before trials, because a missed map makes every other number meaningless.
\item One-shot calibration reads the SLAM canvas, not the room, and never updates.
\end{itemize}

\section{Intellectual property}
This section is a technical briefing for a conversation with a patent attorney. It is not legal advice, not a freedom-to-operate opinion, and not a prior-art search.

\subsection{Facts counsel will need}
\begin{itemize}
\item The repository \texttt{amey1942007/Ros\_lidar\_bot} is public.
\item The explorer's history on that public history begins, for the restored pipeline, at \commit{f493c96} on 12 May 2026. Later public commits add the queue, the flood fill, the phantom counter, and the exact-cell zones through 10 June 2026, and manual zones on 19 July 2026.
\item \texttt{LICENSE.md} is the Apache License, Version 2.0. Section 3 of that licence is a patent grant tied to contributions.
\item No experimental data in the repository establishes an unexpected technical effect.
\item A June 2026 internal report describes an intermediate version and is stored in git history (\commit{29348bf}), on a branch that is also on the public remote.
\item This assessment does not know whether a patent application has already been filed, whether the inventors are employees or students, or who owns the work. Those facts change the answer and have to come from the group and from the institution.
\end{itemize}

\subsection{Why a broad patent is not a realistic goal}
A claim to ``frontier exploration,'' ``scoring frontiers,'' or ``blacklisting failed goals'' will be met with Yamauchi, with utility-based exploration, and with \texttt{explore\_lite}. A claim to flood-filling unreachable space will be met with ordinary occupancy mapping. An examiner does not need this codebase to find those references. A prior-art search is still required before anyone says the residual combination is clear; this note is not that search.

The residual combination (exact-cell enclosed rejection, phantom success as distinct from abort, safety-Boolean preemption into a diverse queue) is narrow. Narrow claims are sometimes granted and often easy to design around: change the threshold, use a disk instead of cell keys, or preempt on Nav2's own stall instead of a custom Boolean. Without a measured effect it is also hard to argue that the combination is more than an aggregation of known recovery tricks.

\subsection{Disclosure and the licence}
Public commits are disclosures. In jurisdictions without a general grace period, a disclosure before filing destroys novelty for what was disclosed. The United States has a one-year grace period for an inventor's own disclosure, and the clock, if it applies, runs from the first public disclosure of each claimed feature, not from the date of this report. India and the European Patent Office should be assumed, until counsel says otherwise, not to offer that grace for a public GitHub history. The authors' commit email domain is an IIT Delhi address. If an Indian filing is the one that matters, the public history is the issue to put in front of counsel first.

Separately, Apache-2.0's patent grant may already license users of this repository to practice the contributions as they appear in the code. Filing later does not, by itself, claw that grant back for those users. Counsel has to read the grant against whatever claim is proposed. The practical consequence is that a patent used to stop people from using this repository is a weak plan even if a patent issues.

Trade secret protection for the algorithm is already gone for everything in the public tree.

Copyright in the code exists and is licensed Apache-2.0. That is a license outward. It is not a substitute for a patent, and it does not reserve an exclusive right to the method.

\subsection{If the group still wants to talk to counsel}
Ask these questions, in this order, and do not file on the strength of this note.
\begin{enumerate}
\item What is the earliest public date of each feature that would be claimed, and has any application already been filed?
\item Who owns the work under the relevant student, employee, or sponsored-research terms?
\item Does Apache-2.0 section 3 already license the proposed claims to users of the repository?
\item After a search, is there any claim that is both narrower than ``frontier exploration'' and still broad enough to be worth the cost?
\item Should any conference submission be delayed until that advice is in hand? Posting a paper is a further disclosure. It is not the first one.
\end{enumerate}
Do not include source code, parameter tables, or this report in a public preprint until those answers exist. The code is already public; the organised description in the two manuscripts is a clearer disclosure than the code alone.

\subsection{What not to claim in a patent draft or a paper}
\begin{itemize}
\item That the policy is optimal, complete, or real-time in a measured sense.
\item That permanent zones are disks of $1.6$\,m. The test is a cell set. The disk is a marker.
\item That exploration terminates.
\item That the June 2026 report's safety cone, phantom threshold, or termination behaviour is the product.
\item Any numerical improvement in coverage, time, or interventions.
\end{itemize}

\section{Recommended order of work}
\begin{enumerate}
\item Decide, with counsel and with the institution, whether a filing is even open. If the answer is no, stop spending time on patent language.
\item Fix the pop score, the border test, and the completion latch. These are correctness issues. They are not new features to patent.
\item Add the per-goal log.
\item Run the protocol in the comprehensive report. Pre-declare the comparisons.
\item Rewrite the concise manuscript around whatever the protocol actually shows, including a negative result.
\item Choose a venue only after that draft exists.
\end{enumerate}

The two manuscripts already written are a starting point for step 5, not a submission. The concise one is the argument. The long one is the specification. Neither should be posted as a preprint until the disclosure question in the previous section is answered and the experimental gap is either filled or stated in the abstract as plainly as it is stated here.

\end{document}
